\scanpage{7}
\mspage{3}
that $\gamma$ is the identity. To see this, one can suppose that $\gamma$ leaves invariant a structure \deleted{of} of an algebraic curve (of anabelian type $(g,\nu)$) on $X$. It suffices \deleted{to} that we have $2g-\nu\geq2$, i.e. that $X$ embeds in \unsure{a|its} generalized Jacobian, extension of an abelian variety by a torus, to be able to conclude from the homological hypothesis that $\gamma$ \deleted{is} induces a \uline{translation} on this Jacobian, by an element of finite order of the homogeneous Jacobian. Thus one must see that $X$ embedded in its Jacobian is not stable under any \deleted{\unsure{automorphism}} nonzero translation, in the anabelian case -- this is a well-known fact\ldots\ [Surely, it should not be difficult to find a purely topological proof of the preceding proposition\ldots]

Let us now note that if $Y$ is an anabelian \inserted{topological} surface, so is every finite étale covering. We thus find the

\deleted{\uline{Proposition}} \inserted{Corollary} 3\quad Let $X,Y$ be \deleted{two} as above, \inserted{$Y$-anabelian}.
Then every isotopy class of finite étale morphisms from $X$ to $Y$ is a right \uline{torsor} under the group $\mathcal{A}_X^0$ of automorphisms of $X$ isotopic to the identity.

In view of proposition 1, it remains
\margin{6}
